Il concetto di query distribuite permette di trarre il meglio dal paradigma del multi-cloud. Una query distribuita è un'interrogazione per l'acquisizione di dati memorizzati su diverse istanze di database. L'utilizzo delle query distribuite permette all'utente di effettuare un'operazioni di selezione, inserimento, modifica e eliminazione su più database contemporaneamente. L'utente che utilizza questo tipo di query è consapevole di interagire con più sorgenti dati ma il risultato che ottiene e la sintassi che utilizza sono analoghi a quelli di una classica interrogazione lanciata su un singolo database. Quest'idea è però un'astrazione di quello che è in realtà il comportamento delle query distribuite. L'esecuzione di una query distribuita su \textit{n} sorgenti non porta all'esecuzione di una singola query in contemporanea su tutte le \textit{n} sorgenti, ma vengono eseguite in successione \textit{n} interrogazioni distinte per ogni singola sorgente. I risultati derivanti da queste \textit{n} interrogazioni vengono poi raggrupati e restituiti all'utente come singolo risultato. \\